\section*{Chapter \ref{ch:b3:bioinformatics} --- Bioinformatics and Sequence Analysis}

\begin{solution}{exo:b3:bioinformatics:1}
Global: both sequences aligned end to end, every residue in a column
--- for two proteins believed homologous over their whole length, such
as \omterm{def:b3:genomics:comparative}{orthologs} of a housekeeping enzyme. Local: the best-scoring pair of
substrings, the rest ignored --- for finding a shared domain (an SH2
domain in two otherwise unrelated signalling proteins), or a gene in a
long genomic sequence.
\end{solution}

\begin{solution}{exo:b3:bioinformatics:2}
Borders $0,-1,-2,-3$ each way. Row A: $1, 0, -1$. Row G: $0, 0, -1$.
Row C: $-1, -1, 1$. Optimum $F(3,3) = 1$: AGC over AAC with no gaps
(match, mismatch, match: $1 - 1 + 1 = 1$).
\end{solution}

\begin{solution}{exo:b3:bioinformatics:3}
$s(a,a) = \lambda^{-1}\log\bigl(q_{aa}/p_{a}^{2}\bigr)$. Tryptophan is
rare ($p_{W} \approx 0.013$), so the chance of two tryptophans aligning
at random, $p_{W}^{2}$, is tiny, and a conserved tryptophan pair is a
far stronger sign of homology than a conserved leucine pair ($p_{L}
\approx 0.1$); the log-odds ratio is correspondingly larger.
\end{solution}

\begin{solution}{exo:b3:bioinformatics:4}
The E-value is the number of \omterm{def:b3:bioinformatics:alignment}{alignments} with a score at least as high
that would be expected by chance in a search of this query against a
database of this size. $E = 3$ means three such scores are expected by
chance: the hit is not evidence of homology (it may still be one, but
the search cannot tell).
\end{solution}

\begin{solution}{exo:b3:bioinformatics:5}
$mn = 400\times 2\times 10^{11} = 8\times 10^{13} = 2^{46.2}$. $E(45) =
2^{1.2} \approx 2.3$; $E(55) = 2^{-8.8} \approx 2\times 10^{-3}$;
$E(65) = 2^{-18.8} \approx 2\times 10^{-6}$. $E = 10^{-3}$ needs $S' =
46.2 + 10.0 = 56$ bits. A \num{40}-residue query has $mn$ ten times
smaller, $2^{42.9}$: $53$ bits suffice --- but a short query can rarely
reach even that.
\end{solution}

\begin{solution}{exo:b3:bioinformatics:6}
\omterm{def:b3:bioinformatics:motif}{Information contents}: $2$, $1$, $0$, and $2 - H$ with $H = -(0.7\log_{2}
0.7 + 3\times 0.1\log_{2} 0.1) = 0.36 + 1.00 = 1.36$, so $0.64$. Total
$R = 3.64$ bits. Chance matches: $9.2\times 10^{6}$ positions on two
strands $\times 2^{-3.64} \approx 7\times 10^{5}$ --- the \omterm{def:b3:bioinformatics:motif}{motif} is
nearly useless alone.
\end{solution}

\begin{solution}{exo:b3:bioinformatics:7}
An insertion of several residues is one mutational event, so its cost
should not grow linearly with its length: an opening cost plus a small
extension cost models this. Too high an opening penalty forces
mismatches where a gap belongs and misaligns everything after a true
insertion; too low a penalty scatters gaps everywhere, matching
residues by chance and inflating identity.
\end{solution}

\begin{solution}{exo:b3:bioinformatics:8}
Not necessarily: the \omterm{def:b3:bioinformatics:alignment}{alignment} covers a \num{130}-residue segment,
which is the signature of a shared domain rather than of an \omterm{def:b3:genomics:comparative}{ortholog}
aligned over its length. Test: search the fly protein back against the
human proteome (is the query its best hit, over the whole length?),
identify the domain with a \omterm{def:b3:bioinformatics:msa}{profile} HMM, and build a gene tree of the
family in several species.
\end{solution}

\begin{solution}{exo:b3:bioinformatics:9}
A \omterm{def:b3:bioinformatics:msa}{profile} records, column by column, what the family tolerates: a
position that is always hydrophobic but never the same residue, an
invariant catalytic residue, a position that is always a gap in half
the family. A pairwise \omterm{def:b3:bioinformatics:alignment}{alignment} scores each residue against one other
residue and cannot know that a valine at position 40 is ``as good as''
the isoleucine there in the query. The \omterm{def:b3:bioinformatics:msa}{profile} also weights the
conserved columns, so weak similarity concentrated where the family is
conserved becomes significant.
\end{solution}

\begin{solution}{exo:b3:bioinformatics:10}
With positive \omterm{prop:b3:bioinformatics:logodds}{expected score} $\mu > 0$ per column, the cumulative score
along the diagonal of two random sequences is a random walk with
positive drift: after $n$ columns it is about $\mu n$, so the best
\omterm{def:b3:bioinformatics:alignment}{local alignment} is essentially the whole thing and its score grows as
$\mu n$ rather than as $\log n$. The Karlin--Altschul theory, which
requires a negative drift so that high scores are rare excursions,
does not apply and no $\lambda$ exists. For DNA at \qty{60}{\%} GC the
chance of a match is $2(0.3^{2}) + 2(0.2^{2}) = 0.26$, so the \omterm{prop:b3:bioinformatics:logodds}{expected
score} is $0.26 - 0.74 = -0.48$: still negative, and the statistics
hold; but a scheme such as match $+1$, mismatch $-0.3$ would have
expectation $+0.04$ and would report the whole genome as one \omterm{def:b3:bioinformatics:alignment}{alignment}.
\end{solution}

\begin{solution}{exo:b3:bioinformatics:11}
\omterm{def:b3:bioinformatics:blast}{Seeds} and chaining: find exact or near-exact matches of $k$-mers between
the two sequences with a hash table, keep the ones that line up on
consistent diagonals, chain them, and run \omterm{thm:b3:bioinformatics:nw}{dynamic programming} only in
the gaps between chained \omterm{def:b3:bioinformatics:blast}{seeds}; it gives up \omterm{def:b3:bioinformatics:alignment}{alignments} in regions with
no \omterm{def:b3:bioinformatics:blast}{seed} (highly divergent stretches). Banding: if the two sequences are
known to be nearly collinear, compute only the cells within a band of
width $w$ around the diagonal, at cost $wn$ instead of $mn$; it gives
up any \omterm{def:b3:bioinformatics:alignment}{alignment} with an insertion larger than the band.
\end{solution}

\begin{solution}{exo:b3:bioinformatics:12}
Coding sequence has a period of three: the three codon positions have
different base compositions (the third is the most biased), and codon
usage is uneven in each species. A model with three coding states in
sequence, each emitting bases with the composition of that codon
position, assigns coding DNA a higher probability than the non-coding
state does, over a window of a few dozen codons, even without stops.
In a human genome exons are short (\qty{150}{bp}), separated by introns
of kilobases, so the coding signal is brief and interrupted; the model
must also recognise splice sites, which are weak signals, and the
sheer amount of non-coding sequence produces many false coding
segments.
\end{solution}

\begin{solution}{pb:b3:bioinformatics:1}
\textbf{1.} Rows K, Q, T; columns K, A, Q, T; borders $0,-1,-2,-3,-4$
and $0,-1,-2,-3$. Row K: $1, 0, -1, -2$; row Q: $0, 0, 1, 0$; row T:
$-1, -1, 0, 2$. Optimum $2$: \texttt{K-QT} over \texttt{KAQT}.
\textbf{2.} Best local score $3$: \texttt{CAT} against \texttt{CAT}
(residues 4--6 of GATCAT with 2--4 of ACAT); \texttt{ATCAT} against
\texttt{A-CAT} also scores $4 - 1 = 3$.
\textbf{3.} $300\times 450 = 1.35\times 10^{5}$ updates; against the
database $300\times 1.2\times 10^{11} = 3.6\times 10^{13}$.
\textbf{4.} $3.6\times 10^{4}$ s, ten hours per query; \omterm{def:b3:bioinformatics:blast}{BLAST}'s \omterm{def:b3:bioinformatics:blast}{seeds}
skip almost all of the table and answer in seconds.
\textbf{5.} $q_{ab} = p_{a}p_{b}e^{\lambda s}$. Alanine: $e^{1.39} = 4.0$,
$q_{AA} = 0.074^{2}\times 4.0 = 0.022$, ratio $4$. Tryptophan:
$e^{3.82} = 45$, $q_{WW} = 0.013^{2}\times 45 = 0.0077$, ratio $45$. An
aligned tryptophan pair is $45$ times more frequent in homologues than
by chance, an alanine pair only four times; alanine pairs are
nevertheless commoner in absolute terms because alanine is common.
\textbf{6.} \qty{24}{\%} lies in the twilight zone, where random
\omterm{def:b3:bioinformatics:alignment}{alignments} reach \qtyrange{15}{20}{\%}; the E-value of the \omterm{def:b3:bioinformatics:alignment}{alignment},
conserved \omterm{def:b3:bioinformatics:motif}{motifs} at the right positions, a profile-HMM match to a known
family, or a shared fold would settle it.
\textbf{7.} $mn = 300\times 1.2\times 10^{11} = 3.6\times 10^{13}$;
$\log_{2}(mn) = 45.0$.
\textbf{8.} $E(92) = 2^{45 - 92} = 2^{-47} \approx 7\times 10^{-15}$;
$E(38) = 2^{7} = 128$.
\textbf{9.} $E = 10^{-3}$ at $S' = 45 + 10 = 55$ bits; $E = 1$ at $45$
bits.
\textbf{10.} Ten times $mn$: $E \approx 7\times 10^{-14}$, still
overwhelming.
\textbf{11.} With $E = 128$ the tenth hit is what chance produces; a
\qty{40}{\%} identity over $60$ residues is not evidence. A \omterm{def:b3:bioinformatics:msa}{profile} search
of residues 200--260 against the domain database could show whether
that segment is a known domain, with statistics that pairwise
comparison lacks.
\textbf{12.} Reciprocal best hits over the full length are consistent
with one-to-one orthology; they do not prove it --- a duplication in
one lineage after the split gives two \omterm{def:b3:genomics:comparative}{co-orthologs}, and the loss of the
true \omterm{def:b3:genomics:comparative}{ortholog} can leave a \omterm{def:b3:genomics:comparative}{paralog} as best hit. A gene tree with several
species is the test.
\textbf{13.} $R = 2 + 2 + 1.6 + 2 + 0.8 + 1.2 + 0.4 + 0.3 = 10.3$ bits.
\textbf{14.} $3.2\times 10^{9}\times 2^{-10.3} \approx 2.5\times 10^{6}$
chance matches.
\textbf{15.} $\log_{2}(3.2\times 10^{9}/200) = \log_{2}(1.6\times 10^{7})
\approx 24$ bits.
\textbf{16.} Co-occurrence at fixed spacing adds the bits: $10.3 + 8 =
18.3$, supplying $8$ of the $13.7$ missing; about $5.7$ bits (a factor
of $50$ in chance matches) must come from elsewhere --- \omterm{def:b3:chromatin-epigenetics:nucleosome}{chromatin}
accessibility, further partners.
\textbf{17.} $H = -(0.5\log_{2}0.5 + 0.5\log_{2}0.5) = 1$ bit; $R = 2 -
1 = 1$ bit.
\textbf{18.} A bacterial factor must find its few sites in a
\qty{4.6}{Mb} genome with no help from \omterm{def:b3:chromatin-epigenetics:nucleosome}{chromatin}: it needs some $19$
bits, and its sites are long and conserved. A eukaryotic genome is a
thousand times larger, needing ten more bits, yet its factors have
short sites; they achieve specificity by combination and by the
restriction of accessible \omterm{def:b3:chromatin-epigenetics:nucleosome}{chromatin}, which also makes regulation more
evolvable, since a short site is easily gained or lost.
\textbf{19.} $3/64 = 0.047$; the expected number of codons before a
stop is $64/3 \approx 21$.
\textbf{20.} $p_{A} = p_{T} = 0.31$, $p_{G} = p_{C} = 0.19$: $P(\text{TAA})
= 0.31^{3} = 0.030$, $P(\text{TAG}) = P(\text{TGA}) = 0.31^{2}\times 0.19
= 0.018$; total $0.066$, expected frame length $15$ codons. AT-rich
DNA is full of stops, so random open frames are shorter and long ones
stand out more.
\textbf{21.} $(61/64)^{100} = e^{-4.80} = 0.008$; $(61/64)^{300} =
e^{-14.4} = 5.6\times 10^{-7}$.
\textbf{22.} Each maximal open frame ends at a stop, and $3.2\times
10^{9}$ codon starts contain $3.2\times 10^{9}\times 3/64 = 1.5\times
10^{8}$ stops: about $1.5\times 10^{8}\times 0.008 = 1.2\times 10^{6}$
random frames of at least $100$ codons, and $1.5\times 10^{8}\times
5.6\times 10^{-7} \approx 80$ of at least $300$.
\textbf{23.} A bacterium of \qty{4.6}{Mb} has some $4\times 10^{5}$ stops
and hence about \num{3500} chance frames of $100$ codons but almost none
of $300$; its genes average $300$ codons and \qty{88}{\%} of the DNA is
coding, so a long open frame is nearly always a gene. In the worm,
\qty{1.5}{\%} of the DNA codes, exons average $50$ codons --- shorter
than the chance threshold --- and a million random frames of $100$
codons swamp them. Eukaryotic gene finders use splice-site signals,
codon bias in a hidden Markov model, homology to known proteins and,
above all, sequenced transcripts.
\textbf{24.} A read from a spliced messenger aligns to the genome in
two pieces separated by an intron: the split marks both splice sites
to the base; read coverage delineates the exons and paired reads link
successive exons into one transcript, resolving which of several
candidate splice sites is used.
\textbf{25.} $E \approx 7\times 10^{-15}$ for the best hit; about $24$
bits to specify $200$ sites in the genome; some $80$ chance reading
frames of $300$ codons in the whole genome.
\end{solution}
